\documentclass[12pt]{article}
\usepackage[utf8]{inputenc}
\usepackage[T1]{fontenc}
\usepackage{hyperref}
\title{Bungu Language Resources for K-12 Curriculum Development}
\author{Compiled by Claude Sonnet 5.5 (claude-sonnet-5-5) for Andr\'as Kornai}
\date{}
\begin{document}
\maketitle

\section{Introduction}
Bungu (Kibungu, Wungu; ISO 639-3 wun; Guthrie F.25) is a Bantu language spoken in the Mbeya region of Tanzania, mainly in Chunya district. Global Recordings Network gives about 36,000 speakers. It should not be confused with Bung, Mbunga, Burunge or Bongo. Little published material turned up in the searches.

\section{Primary Educational Resources}
No primer, textbook or school programme for Bungu turned up in the searches.

\section{Linguistic and Reference Materials}
Wikipedia has a short article: \url{https://en.wikipedia.org/wiki/Bungu_language}. A linguistic data sheet on Bungu numerals, with location details, is at \url{https://lingweb.eva.mpg.de/channumerals/Bungu.htm}. It notes that some numbers use Swahili numerals. Wiktionary has a category page with few entries: \url{https://en.wiktionary.org/wiki/Category:Bungu_language}. The searches turned up no full dictionary or grammar.

\section{Religious and Community Resources}
A Jesus Film in Bungu is reported to exist. I saw no URL for it, so it is not in the JSON. No Bible translation turned up.

\section{Digital and Audio Resources}
The Global Recordings Network language page is \url{https://globalrecordings.net/en/language/wun}. It lists dialect names (Maleza, Mwambani, Songezi, Udinde) and appeared to hold no audio recordings.

\section{Curriculum Development Context}
The language is not known to be taught in schools. Swahili and English are the likely school languages. A mother-tongue curriculum would need new orthography, literacy and reading materials developed with the community.
\end{document}